% The course book's preamble.  Its first part is the union of the preambles
% of the decks the book binds (modules/*/slides*/preamble.tex); the joint
% modules/preamble.tex that the decks will symlink replaces that part (see
% the plan).  Until then this is a copy: *Algoritmiskt tänkande*'s preamble
% plus what the others add (graphicx and the IDE acronym from *Hello,
% World!*, the bibliography URL breaks from *Funktioner*).  The one real
% conflict, cleveref's name for a list item (steg in *Algoritmiskt
% tänkande*, rad in *Hello, World!*), is set per chapter instead: see the
% optional argument of \deck below.
\usepackage[utf8]{inputenc}
\usepackage[T1]{fontenc}
\usepackage[british,english,swedish]{babel}
\usepackage{booktabs}
\usepackage{array}
\usepackage{longtable}

\usepackage[
  natbib,
  style=verbose,
  citestyle=verbose,
  singletitle=false,
  maxbibnames=99,
]{biblatex}
% Merged from the decks' ltnotes.bib by the Makefile.
\addbibresource{ltnotes.bib}
% biblatex breaks a URL only at punctuation, and long blog URLs then overran
% the text block (*Funktioner*).  Allow breaks after digits and at case
% changes as well.
\setcounter{biburlnumpenalty}{7000}
\setcounter{biburlucpenalty}{7000}
\setcounter{biburllcpenalty}{7000}

\usepackage[all]{foreign}
\renewcommand{\foreignfullfont}{}
\renewcommand{\foreignabbrfont}{}

\usepackage{import}
% *Hello, World!* shows screenshots; memoir does not load graphicx itself.
\usepackage{graphicx}

\usepackage[strict]{csquotes}
\usepackage[single]{acro}
% Swedish first, the English term once in parentheses at first use
% (*Hello, World!*).
\DeclareAcronym{IDE}{
  short = IDE,
  long  = {utvecklingsmiljö
           (\foreignlanguage{english}{integrated development environment})},
}

\usepackage{subcaption}

\usepackage[noend]{algpseudocode}
\usepackage{xparse}

\let\email\texttt

% minted v3+ auto-detects the output directory; [outputdir=...] is an error.
% minted prefixes \inputminted's file with the import package's current
% path, so the decks' \inputminted{...}{examples/...} work under \subimport.
\usepackage{minted}
\setminted{autogobble,fontsize=\footnotesize}

% PythonTeX runs from ltxobj; workingdir=.. is book/.  Each chapter moves
% the Python session into its deck's directory (\deck below), so the
% decks' relative paths (examples/..., mentipy.json) resolve as they do
% standalone.
\usepackage[makestderr]{pythontex}
\setpythontexoutputdir{.}
\setpythontexworkingdir{..}

\usepackage{amsmath}
\usepackage{amssymb}
\usepackage{mathtools}
\usepackage{amsthm}
\usepackage{thmtools,thm-restate}
\usepackage[unq]{unique}

\usepackage[binary-units,parse-numbers=false]{siunitx}

\usepackage[colorlinks]{hyperref}
% cleveref reads its language from its own options, not babel's.
\usepackage[swedish]{cleveref}
% memoir and babel call an appendix "bilaga"; cleveref's Swedish names do
% not.  Registered after cleveref's own hook so ours wins.
\AtBeginDocument{%
  \crefname{appendix}{bilaga}{bilagor}%
  \Crefname{appendix}{Bilaga}{Bilagor}%
  \crefname{footnote}{fotnot}{fotnoterna}%
  \Crefname{footnote}{Fotnot}{Fotnoterna}%
}

\usepackage[marginparmargin=outer]{didactic}
% Margin notes and floats share LaTeX's pool of eighteen boxes; the verbose
% citations and \ltnote's empty it without this.
\extrafloats{200}
% A side caption takes part in no collision avoidance; anchoring it at the
% bottom moves it below the \ltnote's.
\ifdefined\setsidecappos\setsidecappos{b}\fi
\ProvideSemanticEnv{lo}{Learning Objective}
  [style=definition,numbered=yes]
  {LO}{LO}
  {Learning objective}{Learning objectives}
\ProvideTranslation{british}{Learning Objective}{Learning Objective}
\ProvideTranslation{british}{LO}{LO}
\ProvideTranslation{british}{Learning objective}{Learning objective}
\ProvideTranslation{british}{Learning objectives}{Learning objectives}
\ProvideTranslation{swedish}{Learning Objective}{Lärandemål}
\ProvideTranslation{swedish}{LO}{lm}
\ProvideTranslation{swedish}{Learning objective}{Lärandemål}
\ProvideTranslation{swedish}{Learning objectives}{Lärandemål}

% Greek letters in the generated hit-list tables of the search appendices.
\DeclareUnicodeCharacter{03C0}{\ensuremath{\pi}}
\DeclareUnicodeCharacter{03BB}{\ensuremath{\lambda}}
\DeclareUnicodeCharacter{03B1}{\ensuremath{\alpha}}
\DeclareUnicodeCharacter{03B2}{\ensuremath{\beta}}
\DeclareUnicodeCharacter{03BC}{\ensuremath{\mu}}


% ------------------------------------------------------------------------
% What the decks' notes.tex adds to the preamble.
% ------------------------------------------------------------------------
\usepackage[noamsthm,notheorems]{beamerarticle}
% The chunks of the woven contents.tex; keep the cross-references, hide
% noweb's ASCII-only identifier index (dbosk/noweb#13).
\usepackage[minted]{noweb}
\noweboptions{breakcode}
\def\nwidentdefs#1{}
\def\nwidentuses#1{}


% ------------------------------------------------------------------------
% Binding the decks.
% ------------------------------------------------------------------------
% A table of contents per chapter.
\usepackage{etoc}

% PythonTeX runs all of a session's code as one program, in document order,
% so a directory change stays in force until the next one.  Each chapter
% changes into its deck's directory: \runpython's examples/..., the stdin
% and deps files, and Mentipy's store and QR directory then resolve as in
% the standalone build.  didactic writes its transcript files in book/
% (DIDACTIC_WORKDIR) whatever the directory, and reads them past the import
% path.  PythonTeX resolves the paths didactic reports to it (created
% files, dependencies) from the book's directory, so they are made
% absolute here; didactic now does so itself, and this override can go
% once the book uses didactic's \runbasedir instead of the directory
% change.
\begin{pycode}
BOOK_DIR = os.getcwd()


def book_enter_deck(path):
    os.chdir(os.path.join(BOOK_DIR, path))


def book_leave_deck():
    os.chdir(BOOK_DIR)


_book_didactic_track = didactic_track


def didactic_track(kind, /, *paths):
    _book_didactic_track(kind, *(os.path.abspath(p) for p in paths))
\end{pycode}

% The decks' appendices (sokprotokoll.tex) are chapters standalone.  In the
% book they follow their chapter, one level down: \chapter becomes
% \section, and so on, and each keeps its letter (A, B, ...), so that a
% pointer such as "föreläsningen Funktioner, bilaga B" holds in the book
% too.  Everything is local to the environment.
\makeatletter
\newenvironment{deckappendices}{%
  \begin{subappendices}%
  \renewcommand{\setthesection}{\Alph{section}}%
  \crefalias{section}{appendix}%
  \let\book@section\section
  \let\book@subsection\subsection
  \let\book@subsubsection\subsubsection
  \let\chapter\book@section
  \let\section\book@subsection
  \let\subsection\book@subsubsection
  \let\subsubsection\paragraph
  % Print the chapter's precis under the section heading, without memoir's
  % pull-up (meant for a chapter heading) and without a ToC entry.
  \let\chapterprecis\chapterprecishere
  \prechapterprecisshift=0pt
}{%
  \end{subappendices}%
}
\makeatother

% Without the decks' in-chapter page breaks (see \deck), the margin notes
% of the densest pages no longer fit beside their markers, and a note
% called in the last lines of a page is carried to the next one.  The book
% sets the margin notes (footnotes, citations and \ltnote's alike:
% didactic's \foottextfont) one size smaller than the decks, and lets the
% margin column run into the bottom margin (marginfix's
% \marginheightadjustment), which is empty: the page number is in the
% running head.
\renewcommand{\foottextfont}{\scriptsize\mpjustification}
\marginheightadjustment=48pt

% \deck[<setup>]{<module>/<deck>}{<title>}: one deck as one chapter.  The
% whole chapter is a biblatex refsection (so the verbose citations start
% afresh and the chapter ends with its own bibliography) and a TeX group
% (so a deck's \newcommand's and <setup> stay in its chapter).  <setup> is
% for the deck-local settings that differ between decks, e.g. cleveref's
% name for a list item.  The deck's appendices are not part of the
% chapter: \deck queues them for \scienceappendices, the book's last part.
\makeatletter
\let\book@appendixchapters\@empty
\NewDocumentCommand{\deck}{O{}mm}{%
  \begin{refsection}%
  \chapter{#3}%
  \label{book:#2}%
  #1%
  \ifcsname c@lo\endcsname\setcounter{lo}{0}\fi
  \pyc{book_enter_deck('../modules/#2')}%
  % The overview, learning objectives and prerequisites first, then the
  % chapter's table of contents at full depth (the book's own lists only
  % parts and chapters), each on its own page(s), and then the chapter's
  % text on a fresh page (author, 2026-09-25).  Starred: no line for it in
  % the book's ToC.
  % Not memoir's abstract environment, as in the decks' notes: its centred
  % title and narrowed text block break the book's flush-left typography.
  % An unnumbered section instead, at the text's full width.
  \section*{\abstractname}
  \subimport{../modules/#2/}{abstract.tex}%
  \clearpage
  \localtableofcontents*
  \clearpage
  % The decks force page breaks inside the chapter (\clearpage in
  % contents.nw and sokprotokoll.tex), each placed to drain the margin-note
  % queue at a page position of the standalone notes.  In the book the
  % pages fall elsewhere, so the breaks only leave short pages.  Within the
  % chapter a \clearpage is a paragraph end; the next chapter's \chapter,
  % outside this group, has the real one.  The margin notes the breaks
  % made room for fit instead in smaller type and a longer margin column
  % (see \foottextfont and \marginheightadjustment above).
  \let\clearpage\par
  \subimport{../modules/#2/}{contents.tex}%
  \printbibliography[heading=subbibintoc]%
  \pyc{book_leave_deck()}%
  \end{refsection}%
  \g@addto@macro\book@appendixchapters{\book@appendixchapter{#1}{#2}{#3}}%
}

% The appendices of every deck, at the end of the book (author,
% 2026-09-27): one chapter per deck, in the order of the decks' chapters,
% named after the deck.  Inside it the deck's appendices keep their
% letters A, B, ... (deckappendices), and since labels are global, the
% chapter text's \cref to them, and pointers such as "föreläsningen
% Funktioner, bilaga B", stay true.  Each appendix chapter is a refsection
% of its own, with the bibliography of the appendices' citations.
\newcommand{\book@appendixchapter}[3]{%
  \begin{refsection}%
  \chapter{#3}%
  #1%
  \pyc{book_enter_deck('../modules/#2')}%
  % A deck with only the method appendix (A) has no searches of its own
  % to point to; the count comes from the previous pass (the .aux).
  \ifnum\book@appendixcount{#2}>1
    Det här kapitlet samlar bilagorna till \cref{book:#2},
    \emph{#3}: hur ett påstående kontrolleras mot
    forskningslitteraturen, och de sökningar som kapitlets påståenden
    vilar på.
  \else
    Det här kapitlet innehåller bilagan till \cref{book:#2},
    \emph{#3}: hur ett påstående kontrolleras mot
    forskningslitteraturen.
  \fi
  \localtableofcontents*
  \let\clearpage\par
  \begin{deckappendices}%
    \subimport{../modules/#2/}{sokprotokoll.tex}%
    \book@recordappendixcount{#2}%
  \end{deckappendices}%
  \printbibliography[heading=subbibintoc]%
  \pyc{book_leave_deck()}%
  \end{refsection}%
}
\newcommand{\scienceappendices}{\book@appendixchapters}
% How many appendices a deck has, as the section counter stands after its
% sokprotokoll.tex inside deckappendices (where each appendix is a
% section).  Written to the .aux and read on the next pass; until then
% (a first pass) a deck counts as having more than the method appendix.
\newcommand{\book@setappendixcount}[2]{%
  \expandafter\xdef\csname book@appcount@#1\endcsname{#2}}
\newcommand{\book@appendixcount}[1]{%
  \ifcsname book@appcount@#1\endcsname
    \csname book@appcount@#1\endcsname
  \else 2\fi}
\newcommand{\book@recordappendixcount}[1]{%
  \immediate\write\@auxout{%
    \string\book@setappendixcount{#1}{\the\value{section}}}}
\makeatother
% The learning objectives restart in every chapter, as standalone, so their
% hyperlink targets must name the chapter.
\AtBeginDocument{\def\theHlo{\theHchapter.\arabic{lo}}}
